\section{Results}\label{sec:results}

\subsection{The Frontend Determines Feasible-Basin Coverage}

Table \ref{tab:frontend-backend} gives the main causal comparison.  With the
frozen atlas, all three backends produce true-feasible and independently
certified recommendations in all 60 domain-seed pairs, with no false
certificate.  Replacing only the frontend by common Sobol reduces these counts
to 0/60 for proposal-only, 1/60 for stacked GP, and 0/60 for canonical SAASBO.
For the preregistered proposal-only and SAAS comparisons, the Holm-adjusted
McNemar $p$-value is $6.94\times10^{-18}$ for both feasibility and
certification.  Figure \ref{fig:frontend-backend} shows the aggregate contrast.

\begin{table}[t]
\centering
\small
\caption{Frontend/backend factorial at $d=1000$.  Frozen-atlas arms use
$n_0=10$ and at most three adaptive calls; proposal-only stops at ten.  Every
arm uses the same independent verifier.}
\label{tab:frontend-backend}
\resizebox{\textwidth}{!}{\begin{tabular}{lllrrrr}
\toprule
Domain & Frontend & Backend & Feasible & Certified & False & Med. regret \\
\midrule
FactorShock & Frozen atlas & Proposal only & 20/20 & 20/20 & 0 & 0.0082 \\
FactorShock & Frozen atlas & Stacked GP & 20/20 & 20/20 & 0 & 0.0082 \\
FactorShock & Frozen atlas & Canonical SAASBO & 20/20 & 20/20 & 0 & 0.0082 \\
FactorShock & Common Sobol & Proposal only & 0/20 & 0/20 & 0 & -- \\
FactorShock & Common Sobol & Stacked GP & 0/20 & 0/20 & 0 & -- \\
FactorShock & Common Sobol & Canonical SAASBO & 0/20 & 0/20 & 0 & -- \\
\midrule
Inventory & Frozen atlas & Proposal only & 20/20 & 20/20 & 0 & 0.0116 \\
Inventory & Frozen atlas & Stacked GP & 20/20 & 20/20 & 0 & 0.0116 \\
Inventory & Frozen atlas & Canonical SAASBO & 20/20 & 20/20 & 0 & 0.0031 \\
Inventory & Common Sobol & Proposal only & 0/20 & 0/20 & 0 & -- \\
Inventory & Common Sobol & Stacked GP & 1/20 & 1/20 & 0 & 0.0108 \\
Inventory & Common Sobol & Canonical SAASBO & 0/20 & 0/20 & 0 & -- \\
\midrule
Queue & Frozen atlas & Proposal only & 20/20 & 20/20 & 0 & 0.0029 \\
Queue & Frozen atlas & Stacked GP & 20/20 & 20/20 & 0 & 0.0029 \\
Queue & Frozen atlas & Canonical SAASBO & 20/20 & 20/20 & 0 & 0.0015 \\
Queue & Common Sobol & Proposal only & 0/20 & 0/20 & 0 & -- \\
Queue & Common Sobol & Stacked GP & 0/20 & 0/20 & 0 & -- \\
Queue & Common Sobol & Canonical SAASBO & 0/20 & 0/20 & 0 & -- \\
\bottomrule
\end{tabular}
}
\end{table}

\begin{figure}[t]
\centering
\includegraphics[width=0.82\textwidth]{frontend_backend.pdf}
\caption{Feasible-basin coverage is driven by the frozen structural frontend,
not by the backend.  Counts aggregate 20 seeds in each of three domains.}
\label{fig:frontend-backend}
\end{figure}

Once coverage is established, the backend affects objective quality.  Across
all 60 frozen-atlas pairs, canonical SAASBO improves regret over proposal-only
in 29 seeds, loses in two, and ties in 29; the Holm-adjusted Wilcoxon
$p$-value is $1.16\times10^{-6}$.  It improves over stacked GP in 29 seeds,
loses in three, and ties in 28 ($p_{\mathrm{Holm}}=1.49\times10^{-6}$).
The median paired difference is zero because FactorShock and many other pairs
tie, while Inventory and Queue contain the improvements.  The target-stage
convergence paths in Figure \ref{fig:convergence} make the decomposition clear:
the initial atlas reaches the feasible region, and the three adaptive calls
occasionally improve its objective.

\begin{figure}[t]
\centering
\includegraphics[width=0.86\textwidth]{target_convergence.pdf}
\caption{Median best true-feasible objective gap during target search.  Missing
early values mean that no evaluated point was truly feasible at that call.
Verification samples are excluded.}
\label{fig:convergence}
\end{figure}

\subsection{Both Atlas Components Matter}

Universal structural support alone reaches 27/60 feasible trials; source
templates alone reach 40/60; their combined maximin atlas reaches 60/60
(Table \ref{tab:components} and Figure \ref{fig:components}).  Universal
support is particularly important as a source-misspecification fallback,
whereas source templates concentrate probability mass around empirically safe,
low-cost profiles.  The combined atlas is therefore not merely the best source
profile copied to a new dimension.

\begin{table}[t]
\centering
\caption{Causal frontend component ablation with the same source archive,
target seeds, proposal budget, and verifier.}
\label{tab:components}
\begin{tabular}{lrrrr}
\toprule
Frontend component & Feasible & Certified & False & Med. regret \\
\midrule
Universal support & 27/60 & 23/60 & 0 & 0.0082 \\
Source templates & 40/60 & 40/60 & 0 & 0.0072 \\
Combined maximin atlas & 60/60 & 60/60 & 0 & 0.0082 \\
\bottomrule
\end{tabular}

\end{table}

\begin{figure}[t]
\centering
\includegraphics[width=0.70\textwidth]{frontend_components.pdf}
\caption{Feasible and independently certified recommendations under the
frontend component ablation.}
\label{fig:components}
\end{figure}

\subsection{Archive-Fair Transfer and Total-Cost Controls}

Table \ref{tab:transfer} compares transfer methods under exactly the same
source archive, frozen initial design, target budget, seeds, and verifier.
Because every method receives the strong atlas, most achieve high feasible
coverage.  Safe F-PACOH reaches 56/60, RGPE and FSBO 59/60, transfer MTGP
57/60, and stacked GP, HyperBO, MetaBO, MALIBO, and canonical SAASBO 60/60.
Canonical SAASBO has the lowest aggregate median feasible regret (0.0038),
whereas the other displayed methods have 0.0082.  Verification cost differs
materially, demonstrating why the 60/60 coverage count alone is insufficient.

\begin{table}[t]
\centering
\small
\caption{Archive-fair transfer comparison at $d=1000$, $S=384$, $n_0=10$,
and $N=13$.  Mean verification includes safety and any fresh objective-
comparison calls.}
\label{tab:transfer}
\begin{tabular}{lrrrrr}
\toprule
Method & Feasible & Certified & False & Med. regret & Mean verify \\
\midrule
Atlas + canonical SAASBO & 60/60 & 60/60 & 0 & 0.0038 & 173.9 \\
Safe F-PACOH & 56/60 & 56/60 & 0 & 0.0082 & 215.5 \\
RGPE & 59/60 & 59/60 & 0 & 0.0082 & 154.9 \\
Stacked GP & 60/60 & 60/60 & 0 & 0.0082 & 90.9 \\
Transfer MTGP & 57/60 & 57/60 & 0 & 0.0082 & 204.3 \\
FSBO & 59/60 & 59/60 & 0 & 0.0082 & 181.6 \\
HyperBO & 60/60 & 60/60 & 0 & 0.0082 & 211.7 \\
MetaBO & 60/60 & 60/60 & 0 & 0.0082 & 158.9 \\
MALIBO & 60/60 & 60/60 & 0 & 0.0082 & 162.9 \\
\bottomrule
\end{tabular}

\end{table}

The equal-total-search-cost comparison is more severe.  Target-only TuRBO and
SCBO receive 397 target-search calls and find no feasible recommendation in 60
trials.  Periodic-capped target-only SAAS finds two, both certified.  The atlas
plus canonical SAASBO uses 384 source calls and 13 target-search calls and
certifies all 60 (Table \ref{tab:totalcost}).  The Holm-adjusted McNemar
$p$-values for the final method versus each target-only control are below
$3\times10^{-17}$.  These results support reuse of related source simulations;
they do not imply that source calls are free or that the frontend is uniformly
better than every task-specific structured design.

\begin{table}[t]
\centering
\caption{Equal source-plus-search-cost comparison.  Verification is applied
after search and reported separately.  ``Target-only SAASBO'' denotes the
audited periodic-capped control, not canonical every-iteration SAASBO.}
\label{tab:totalcost}
\begin{tabular}{lrrrrr}
\toprule
Method & Source & Search & Feasible & Certified & False \\
\midrule
Atlas + canonical SAASBO & 384 & 13 & 60/60 & 60/60 & 0 \\
Target-only TuRBO & 0 & 397 & 0/60 & 0/60 & 0 \\
Target-only SCBO & 0 & 397 & 0/60 & 0/60 & 0 \\
Target-only SAASBO & 0 & 397 & 2/60 & 2/60 & 0 \\
\bottomrule
\end{tabular}

\end{table}

\subsection{Dimension and Target-Budget Frontier}

At $d=1000$, proposal-only at ten search calls and canonical SAASBO at 13 both
certify all 60 trials; aggregate median regret improves from 0.00825 to 0.00383.
At $d=10000$, all 30 trials are feasible and certified with zero false
certificates at search budgets 10, 13, 20, and 40 (Table
\ref{tab:dimension}).  The corresponding $d/N$ ratios range from 1,000 to 250.
Median regret is 0.00825, 0.00825, 0.00451, and 0.00579.  The deterioration
from $N=20$ to $N=40$ is retained rather than smoothed away; additional
adaptive search is not monotonically beneficial under this finite stochastic
protocol.

\begin{table}[t]
\centering
\small
\caption{Dimension and target-search budget.  Source archive size remains
384; verification is additional and variable.}
\label{tab:dimension}
\begin{tabular}{rrrrrrrr}
\toprule
$d$ & Search & $d/N$ & Source & Feasible & Certified & False & Med. regret \\
\midrule
1000 & 10 & 100.0 & 384 & 60/60 & 60/60 & 0 & 0.0082 \\
1000 & 13 & 76.9 & 384 & 60/60 & 60/60 & 0 & 0.0038 \\
10000 & 10 & 1000.0 & 384 & 30/30 & 30/30 & 0 & 0.0082 \\
10000 & 13 & 769.2 & 384 & 30/30 & 30/30 & 0 & 0.0082 \\
10000 & 20 & 500.0 & 384 & 30/30 & 30/30 & 0 & 0.0045 \\
10000 & 40 & 250.0 & 384 & 30/30 & 30/30 & 0 & 0.0058 \\
\bottomrule
\end{tabular}

\end{table}

\begin{figure}[t]
\centering
\includegraphics[width=0.90\textwidth]{dimension_budget.pdf}
\caption{Dimension/budget frontier: aggregate median feasible regret (left)
and mean verification calls (right).  The $d=10000$ result is an effective-
dimension stress test under the registered structural family.}
\label{fig:dimension}
\end{figure}

The result is consistent with Corollary \ref{cor:dimension}: the raw dimension
can grow while the profile coordinate remains fixed.  It is not evidence that
SAASBO or the atlas can identify an arbitrary sparse set of raw coordinates in
ten calls.  Indeed, the common-Sobol negative controls show that the same
backend without structural coverage fails at $d=1000$.

\subsection{External Electricity Holdout}\label{sec:energy-results}

All three energy arms are independently certified in 20/20 seeds with no false
certificate (Table \ref{tab:energy}).  The frozen atlas at $N=13$ beats
equal-total-cost common Sobol at $N=397$ in all 20 paired seeds.  The median
objective difference (atlas minus Sobol) is $-0.05331$, with bootstrap 95\%
interval $[-0.05408,-0.05058]$ and one-sided exact sign
$p=9.54\times10^{-7}$.

The natural low-frequency grid changes the interpretation.  The atlas wins
9/20 and loses 11/20; the median difference is $0.000945$, with 95\% interval
$[-0.00108,0.00180]$ and one-sided sign $p=0.748$.  Thus the external result
supports low-frequency structural compression and a large advantage over
unstructured total-cost search.  It does not establish that source learning
outperforms an obvious domain-specific low-frequency control.

\begin{table}[t]
\centering
\caption{Untouched Great Britain electricity-storage holdout and retained
post-confirmatory fairness controls.}
\label{tab:energy}
\begin{tabular}{lrrrr}
\toprule
Arm & Certified & False & Median objective & Mean verify \\
\midrule
Frozen atlas, $N=13$ & 20/20 & 0 & 0.0976 & 80.0 \\
Natural low-frequency grid, $N=13$ & 20/20 & 0 & 0.0956 & 80.0 \\
Common Sobol, $N=397$ & 20/20 & 0 & 0.1488 & 80.0 \\
\bottomrule
\end{tabular}

\end{table}

\begin{figure}[t]
\centering
\includegraphics[width=0.90\textwidth]{external_energy.pdf}
\caption{External energy paired differences.  Negative values favor the
frozen atlas.  The left panel is a retained negative result; the right panel
is the equal-total-cost unstructured control.}
\label{fig:energy}
\end{figure}

\subsection{Heteroscedastic Decomposition Is a Calibration Diagnostic}
\label{sec:hvd-results}

The cumulative factor model sharply improves variance-shape recovery under a
matched proposal, canonical SCBO backend, verifier, and 20-seed design.  Mean
log-variance RMSE falls from 1.158 to 0.219 in FactorShock, 0.577 to 0.349 in
Inventory, and 0.501 to 0.254 in Queue.  Shape correlation rises from zero
under pooled variance to 0.919, 0.953, and 0.995.  The Holm-adjusted Wilcoxon
$p$-value is $1.14\times10^{-10}$ for both log-RMSE and shape correlation.

Operationally, however, both models produce 59/60 feasible and certified
recommendations, zero false certificates, and identical median paired regret.
The HVD model also increases mean verification cost in Inventory and Queue.
It therefore fails the preregistered promotion rule and remains an optional
risk-calibration diagnostic.

\begin{table}[t]
\centering
\small
\caption{Pooled versus cumulative factor-HVD.  Variance diagnostics are
post-run only; they do not enter target selection.}
\label{tab:hvd}
\begin{tabular}{lllrrrr}
\toprule
Domain & Variance model & Feasible & Log-RMSE & Shape corr. & Mean verify \\
\midrule
FactorShock & Pooled & 20/20 & 1.158 & 0.000 & 159.2 \\
FactorShock & Cumulative factor-HVD & 20/20 & 0.219 & 0.919 & 158.4 \\
\midrule
Inventory & Pooled & 20/20 & 0.577 & 0.000 & 182.4 \\
Inventory & Cumulative factor-HVD & 20/20 & 0.349 & 0.953 & 184.0 \\
\midrule
Queue & Pooled & 19/20 & 0.501 & 0.000 & 190.4 \\
Queue & Cumulative factor-HVD & 19/20 & 0.254 & 0.995 & 197.6 \\
\bottomrule
\end{tabular}

\end{table}

\begin{figure}[t]
\centering
\includegraphics[width=0.78\textwidth]{hvd_diagnostic.pdf}
\caption{Cumulative factor-HVD recovers variance shape but does not improve
feasible recovery or regret in the matched causal experiment.}
\label{fig:hvd}
\end{figure}
